\chapter{Documento de Arquitetura de Software (SAD)}
\label{chap:arquitetura_software}

\section{Abordagem de Descrição Arquitetural}
A descrição arquitetural do sistema segue os preceitos da norma internacional \textbf{ISO/IEC/IEEE 42010:2022} e o framework de completude \textbf{ArchCaMo} (\textit{Architecture Completeness Model}). A arquitetura é estruturada em visões arquiteturais baseadas no modelo 4+1 de Kruchten, complementadas por uma visão dedicada de Segurança e Conformidade (LGPD).

\section{Sistema de Interesse, Ambiente e Escopo Arquitetural}
\begin{itemize}
    \item \textbf{Sistema de Interesse:} Plataforma integrada orientada a serviços com persistência relacional e interface web responsiva;
    \item \textbf{Ambiente Operacional:} Infraestrutura em nuvem (IaaS/PaaS) em contêineres Docker orquestrados;
    \item \textbf{Limites de Escopo:} Não inclui integração com hardware médico proprietário de aquisição de sinais vitais na versão atual.
\end{itemize}

\section{Stakeholders e Concerns Arquiteturais}
\subsection{Stakeholders Arquiteturais}
\begin{table}[htbp]
\caption{Mapeamento dos Stakeholders Arquiteturais}
\label{tab:stakeholders_arquitetura}
\centering
\small
\begin{tabularx}{\textwidth}{@{} l l Y @{}}
\toprule
\textbf{ID} & \textbf{Stakeholder} & \textbf{Interesse Principal na Arquitetura} \\
\midrule
S1 & Desenvolvedores Front-end & Contratos de API REST claros, tipagem consistente e baixa latência. \\
S2 & Desenvolvedores Back-end & Separação limpa de camadas, isolamento de regras de negócio e ORM assíncrono. \\
S3 & Engenheiro de DevOps/SRE & Facilidade de deploy contínuo, observabilidade, logs e contenção em Docker. \\
S4 & Encarregado de Dados (DPO) & Rastreabilidade de acessos a dados sensíveis, conformidade estrita com a LGPD. \\
S5 & Gestor de Negócio & Disponibilidade contínua da operação, baixo custo operacional e escalabilidade. \\
\bottomrule
\end{tabularx}
\end{table}

\subsection{Concerns Arquiteturais (Preocupações)}
\begin{table}[htbp]
\caption{Concerns Arquiteturais Catalogados}
\label{tab:concerns_arquitetura}
\centering
\small
\begin{tabularx}{\textwidth}{@{} l l Y @{}}
\toprule
\textbf{ID} & \textbf{Concern} & \textbf{Descrição do Requisito de Arquitetura} \\
\midrule
C1 & Segurança e Autenticação & Garantir isolamento de privilégios (RBAC) e autenticação segura stateless via JWT. \\
C2 & Desempenho e Concorrência & Evitar colisões transacionais em reservas simultâneas via transações atômicas ACID. \\
C3 & Manutenibilidade e Testabilidade & Baixo acoplamento entre camadas para viabilizar cobertura de testes $\ge 80\%$. \\
C4 & Resiliência e Recuperação & Mecanismos de backup automatizado diário e procedimento de rollback sem perda de dados. \\
\bottomrule
\end{tabularx}
\end{table}

\section{Viewpoints e Views Arquiteturais}

\subsection{Visão de Cenários (Casos de Uso Arquiteturalmente Significantes)}
Cenários críticos que exercitam as propriedades não-funcionais do sistema (ex.: alta concorrência e operações financeiras).

\subsection{Visão Lógica}
A Visão Lógica descreve a organização em componentes e subsistemas independentes. A divisão adotada segue princípios de Contextos Delimitados (Bounded Contexts) para garantir coesão interna e baixo acoplamento.

\subsection{Visão de Processos}
Modela o fluxo de execução, transações assíncronas, controle de concorrência e comunicações entre serviços.

\subsection{Visão de Desenvolvimento}
Estruturação do código-fonte em camadas limpas (ex.: Camada de Apresentação, Aplicação/Casos de Uso, Domínio e Infraestrutura/Persistência).

\subsection{Visão Física / de Implantação}
Mapeamento dos nós computacionais, balanceadores de carga, servidores de aplicação em contêineres e instâncias gerenciadas de banco de dados.

\subsection{Visão de Segurança e Conformidade (LGPD)}
Todos os dados pessoais sensíveis (prontuários, diagnósticos e documentos) são criptografados em repouso (AES-256) e em trânsito (TLS 1.3). Cada leitura ou modificação gera log imutável de auditoria com identificação do operador, data/hora e IP.

\section{Decisões Arquiteturais e Rationale (ADRs)}
As decisões arquiteturais fundamentais do sistema são formalizadas a seguir:

\subsection{ADR-01: Adoção de Arquitetura Modular em Camadas}
\textbf{Status:} Aprovado \quad|\quad \textbf{Data:} 2026-09-14

\paragraph{Contexto:} Necessidade de viabilizar desenvolvimento paralelo por equipe distribuída e facilidade de testes unitários isolados.
\paragraph{Decisão:} Implementar padrão em camadas (Presentation, Application, Domain, Infrastructure).
\paragraph{Consequências:} Redução de acoplamento, maior manutenibilidade e total independência do framework web em relação às regras de domínio.

\subsection{ADR-02: Banco de Dados Relacional PostgreSQL com ORM Assíncrono}
\textbf{Status:} Aprovado \quad|\quad \textbf{Data:} 2026-09-14

\paragraph{Contexto:} Operações transacionais críticas (estoque, financeiro e agendamento) exigem integridade ACID absoluta.
\paragraph{Decisão:} Utilizar PostgreSQL 16 com suporte a migrações versionadas de esquema.
\paragraph{Consequências:} Garantia formal de integridade referencial e suporte a alto volume de leitura/escrita simultânea.

\section{Matriz de Rastreabilidade (Concerns vs Decisões vs Views)}
\begin{table}[htbp]
\caption{Matriz de Rastreabilidade Arquitetural}
\label{tab:rastreabilidade_arquitetura}
\centering
\small
\begin{tabularx}{\textwidth}{@{} l l l Y @{}}
\toprule
\textbf{Concern} & \textbf{Visão Associada} & \textbf{Decisão (ADR)} & \textbf{Mecanismo de Validação} \\
\midrule
C1 (Segurança) & Segurança / Lógica & ADR-01 (Camadas) & Testes de autorização RBAC e verificação de token. \\
C2 (Concorrência) & Processos / Física & ADR-02 (PostgreSQL) & Teste de estresse com transações simultâneas no agendamento. \\
C3 (Manutenibilidade) & Desenvolvimento & ADR-01 (Camadas) & Cobertura de código e análise estática (linting/sonarqube). \\
\bottomrule
\end{tabularx}
\end{table}
